Our construction will make use of a protocol between a client and a server that allows the client to learn a point on a pseudorandomly derived polynomial such that the polynomial is hidden from \emph{both} parties.
We say that such a protocol realizes \emph{double-oblivious polynomial evaluation} (DOPE), and give an ideal functionality for it in
\cref{fig:ideal_dope}.

\begin{figure}[h]
	\begin{center}
		\begin{techbox}
			\centerline{\textbf{Functionality} $\fdope$}

			This functionality is parameterized by a field $\field$ and an integer $\ell$. For a field $\field$, integer $\ell$, and field element $\val \in \field$, let $\Polys_{\field, \ell, \val}$ be the set of all polynomials over $\field$ of degree $\ell$ with leading coefficient $1$ where the constant term is $\val$.

			\boxsep

			Ignore all messages associated with $\sid$ before receiving $(\init, \sid)$ from $\randserver$.


			\begin{algenumerate}
				\item Upon receiving $(\init, \sid)$ from $\randserver$, initialize an empty table $T$ and store $T$.
				\item Upon receiving $(\eval, \sid, \val, x)$ from a user $\user$, send $(\eval, \sid, \user)$ to $\simulator$.
				Then, if $\val \notin T$, sample $p \sample \Polys_{\field, \ell, \val}$ and set $T[\val] := p$.
				Finally, retrieve $p \coloneqq T[v]$, and return $p(x)$ to $\user$.

				\item Upon receiving $(\observe, \sid, \val)$ from $\simulator$, ignore the message if $\randserver$ is not corrupt.
				Else, retrieve $T$ and if $\val \notin T$, sample $p \sample \Polys_{\field, \ell, \val}$ and set $T[\val] \coloneqq p$.
				Finally, retrieve $p \coloneqq T[v]$, and return $p$ to $\simulator$.
			\end{algenumerate}
		\end{techbox}
	\end{center}
	\caption{Ideal functionality for DOPE.}
	\label{fig:ideal_dope}
\end{figure}


We describe a protocol realizing the functionality in \cref{fig:protocol_dope}.
Our protocol uses a linearly homomorphic public-key encryption scheme with a dense ciphertext space.

Several encryption schemes---such as Paillier~\autocite{EC:Paillier99}, Damg\aa{}rd-Jurik~\autocite{PKC:DamJur01}, Goldwasser-Micali~\autocite{STOC:GolMic82}, and Benaloh~\autocite{SAC:Ben94}---are linearly homomorphic and have dense ciphertext space.
In this work, we use the Benaloh encryption scheme, described in \cref{sec:benaloh_scheme}.
We use this scheme as it supports having a polynomial order finite field as the message space, which is what is will be required by the MDSS scheme used in our $t$-heavy-hitters construction.

\begin{figure}[h]
	\begin{center}
		\begin{techbox}
			\centerline{\textbf{Protocol} $\pdope[\ell]$}
			Let $\field$ be a finite field.
			Let $\encscheme$ be a linearly homomorphic (\cref{def:enc-linear-homomorphism}) public-key encryption scheme with dense ciphertext space (\cref{def:dense-ctx-space}), and with message space $\field$.

			\boxsep
			When a user $\user$ is activated with input $(\eval, \sid, \val, x)$, it proceeds as follows:
			\begin{algenumerate}
				\item $\user$ invokes $\fsetup$ to receive $\pk$.
				Then, for $i \in \nrange{\ell}$, it sends $(\query, \sid, \val \concat i)$ to $\frandoracle$ and receives $\ct_i$.

				\item $\user$ samples $\dopemask \gets \field$ uniformly at random and computes $\ct_0 \gets \encscheme.\Enc(\pk, v + \dopemask)$.

				\item $\user$ computes
				\begin{equation*}
					\dopectmasked \coloneqq \encscheme.\Eval\parens{\pk, (\ct_0, \dots, \ct_{\ell}), (1, x, \dots, x^{\ell})}
				\end{equation*}
				and sends $\dopectmasked$ to $\randserver$.

				\item Upon receiving $\dopectmasked$ from $\user$, $\randserver$ invokes $\fsetup$ to receive $(\pk, \sk)$, computes $\dopemaskedshare \coloneqq \encscheme.\Dec(\sk, \dopectmasked)$, and sends $\dopemaskedshare$ to $\user$.

				\item Upon receiving $\dopemaskedshare$ from $\randserver$, $\user$ computes $\dopeoutput \coloneqq \dopemaskedshare - \dopemask$ and outputs $\dopeoutput$.
			\end{algenumerate}
		\end{techbox}
	\end{center}
	\caption{Protocol for Doubly-Oblivious Polynomial Evaluation.}
	\label{fig:protocol_dope}
\end{figure}

